\section[Learning a classifier]{Learning a classifier: likelihood, loss and training}
\label{sec:bg:learning}

Every learned model of the thesis, trees and networks alike, is trained in the same way:
it is a function with parameters that turns the description of a sample into three
probabilities, and the parameters are chosen to make the probabilities of the observed
labels as large as possible on the training samples.  This section spells that out, in
the order in which it happens: the model as a conditional distribution, the likelihood and
the loss that comes from it, the gradient and the optimiser that follows it, the
corrections for class imbalance, and the guards against overfitting.  \Cref{fig:bg:loop} is the loop that the
rest of the section describes.

\begin{figure}[t]
\centering
\resizebox{\textwidth}{!}{\begin{tikzpicture}[>=Latex, x=1cm, y=1cm,
  box/.style={draw=vizInk, rounded corners=2pt, align=center, font=\footnotesize, inner sep=4pt, minimum height=10mm, text width=27mm, fill=white},
  input/.style={box, dashed}, op/.style={box, fill=vizB!12}, state/.style={box, fill=vizA!12}, outbox/.style={box, fill=vizC!15},
  arr/.style={->, draw=vizInk, line width=0.8pt}, lab/.style={font=\scriptsize, text=vizInk, align=center}]
\node[input] (data)  at (0,0)     {mini-batch of samples\\$(\mathbf{x}_i, y_i)$, $i \in B$};
\node[state] (theta) at (0,-2.8)  {parameters $\boldsymbol{\theta}$\\(weights and biases)};
\node[op]    (fwd)   at (3.8,0)   {forward pass\\$\hat{\mathbf{p}}_i = f(\mathbf{x}_i; \boldsymbol{\theta})$};
\node[op]    (loss)  at (7.6,0)   {loss\\$L = \frac{1}{|B|}\sum_i \ell(\hat{\mathbf{p}}_i, y_i)$};
\node[op]    (bwd)   at (11.4,0)  {backward pass\\$\nabla_{\boldsymbol{\theta}} L$ by the chain rule};
\node[op]    (upd)   at (11.4,-2.8) {update\\$\boldsymbol{\theta} \leftarrow \boldsymbol{\theta} - \eta\, \tilde{\nabla}_{\boldsymbol{\theta}} L$};
\node[outbox] (val)  at (5.7,-5.0) {after every epoch: score on the validation day;\\stop when it stops improving};
\draw[arr] (data) -- (fwd);
\draw[arr] (theta) -- (fwd);
\draw[arr] (fwd) -- (loss);
\draw[arr] (loss) -- (bwd);
\draw[arr] (bwd) -- (upd);
\draw[arr] (upd) -- node[lab, above] {new $\boldsymbol{\theta}$, next mini-batch} (theta);
\draw[arr, dashed] (upd.south) |- (val.east);
\draw[arr, dashed] (val.west) -| (theta.south);
\end{tikzpicture}
}
\caption{The training loop of a neural classifier.  A mini-batch of samples goes through
the model (forward pass), the loss measures how far the predicted probabilities are from
the labels, backpropagation computes the gradient of the loss with respect to every
parameter (backward pass), and the optimiser moves the parameters a small step against the
gradient; after every pass over the training data the model is scored on the validation
day, and training stops when that score stops improving.}
\label{fig:bg:loop}
\end{figure}

\subsection{The model as a conditional distribution}
\label{sec:bg:learning:model}

A classifier with three classes is a function $f(\mathbf{x}; \boldsymbol{\theta})$ that
maps the description $\mathbf{x}$ of a sample to a vector of three probabilities
$\hat{\mathbf{p}} = (\hat{p}_1, \hat{p}_2, \hat{p}_3)$, one per class, that sum to one.
Inside, the function first computes three unconstrained numbers, the \emph{logits}
$\mathbf{z} = (z_1, z_2, z_3)$, and then turns them into probabilities with the
\emph{softmax} function,
\begin{equation}
  \hat{p}_c = \operatorname{softmax}(\mathbf{z})_c = \frac{e^{z_c}}{\sum_{c'} e^{z_{c'}}},
  \qquad c = 1, 2, 3,
  \label{eq:softmax}
\end{equation}
which makes every $\hat{p}_c$ positive and the three sum to one, and which preserves the
order of the logits: the largest logit is the most probable class.  Adding the same
constant to every logit changes nothing; adding a constant to one logit raises that
class's probability at the expense of the others, which is what the calibration offsets
$\boldsymbol{\delta}$ of \cref{sec:bg:notation} do.  The model is therefore a
\emph{conditional probability distribution} of the label given the inputs,
$p(y = c \mid \mathbf{x}; \boldsymbol{\theta}) = \hat{p}_c$, and training is the choice of
$\boldsymbol{\theta}$.

\subsection{Maximum likelihood and the cross-entropy loss}
\label{sec:bg:learning:likelihood}

Given $N$ training samples $(\mathbf{x}_i, y_i)$ that are treated as independent, the
probability that the model assigns to the observed labels, the \emph{likelihood}, is the
product of the probabilities it gives to each true class,
\begin{equation}
  \mathcal{L}(\boldsymbol{\theta}) = \prod_{i=1}^{N} p(y_i \mid \mathbf{x}_i; \boldsymbol{\theta})
  = \prod_{i=1}^{N} \hat{p}_{i, y_i},
\end{equation}
where $\hat{p}_{i, y_i}$ is the probability that the model gives, for sample $i$, to the
class $y_i$ that actually occurred.  The \emph{maximum-likelihood} principle chooses the
parameters that make this product largest \cite{bishop2006prml}.  A product of millions
of numbers below one is unusable as it stands, so one takes its logarithm, which turns the
product into a sum and does not change where the maximum is, and one changes its sign, so
that the quantity is to be minimised and is a \emph{loss}:
\begin{equation}
  L(\boldsymbol{\theta}) = -\frac{1}{N} \sum_{i=1}^{N} \log \hat{p}_{i, y_i}
  = \frac{1}{N} \sum_{i=1}^{N} \ell(\hat{\mathbf{p}}_i, y_i),
  \qquad
  \ell(\hat{\mathbf{p}}, y) = -\log \hat{p}_{y} = -\sum_{c} [y = c] \log \hat{p}_c .
  \label{eq:xent}
\end{equation}
This is the \emph{cross-entropy} loss: it is zero when the model gives probability one to
the true class and grows without bound as that probability goes to zero, so a confident
mistake is punished far more than a hesitant one.  The $1/N$ makes the loss an average
per sample.  Minimising the cross-entropy is the same as maximising the likelihood, and it
has a second reading: it is the average number of nats of surprise that the labels cause
the model, which is why a model with a lower cross-entropy is said to ``fit'' the data
better.  The cross-entropy is also what makes the outputs of the model \emph{probabilities}
in the useful sense: among all functions, the one that minimises it on infinitely many
samples is the true conditional distribution $p(y \mid \mathbf{x})$, a property called
\emph{properness} that the squared error, say, does not have for classification.

With the softmax of \cref{eq:softmax} the derivative of this loss with respect to the
logits is remarkably simple,
\begin{equation}
  \frac{\partial \ell}{\partial z_c} = \hat{p}_c - [y = c],
  \qquad \text{or in vector form} \qquad
  \frac{\partial \ell}{\partial \mathbf{z}} = \hat{\mathbf{p}} - \mathbf{e}_{y},
  \label{eq:xent-grad}
\end{equation}
the predicted probabilities minus the one-hot vector of the truth: the gradient pushes
the logit of the true class up by the amount the model under-estimated it and the logits of
the other classes down by the amount it over-estimated them.  Every gradient computed in
the thesis starts from this vector.

\subsection{Gradient descent and its stochastic version}
\label{sec:bg:learning:sgd}

The loss is a function of millions of parameters and cannot be minimised in closed form.
\emph{Gradient descent} starts from random parameters and repeats a small step in the
direction in which the loss decreases fastest, the negative of its gradient:
\begin{equation}
  \boldsymbol{\theta} \leftarrow \boldsymbol{\theta} - \eta\, \nabla_{\boldsymbol{\theta}} L(\boldsymbol{\theta}),
  \label{eq:gd}
\end{equation}
where $\nabla_{\boldsymbol{\theta}} L$ is the vector of the partial derivatives of the loss
with respect to every parameter and $\eta$, the \emph{learning rate}, is the size of the
step.  Computing the gradient of the full loss \cref{eq:xent} would require a pass over
all $N$ samples for every step, which is wasteful: the average over a random subset, a
\emph{mini-batch} $B$ of a few thousand samples, is an unbiased estimate of the gradient
and costs a thousandth of the time.  Replacing the full gradient with the mini-batch
gradient gives \emph{stochastic gradient descent}, and one pass of mini-batches over the
whole training set is an \emph{epoch}.  The noise of the estimate is not only harmless but
useful: it lets the parameters escape the flat regions and poor minima in which the exact
gradient would get stuck.

Plain stochastic gradient descent uses the same learning rate for every parameter, a poor
choice when the parameters live on very different scales, as they do in a deep network.
\emph{Adam} \cite{kingma2015adam}, in its AdamW variant \cite{loshchilov2019adamw}, is
the optimiser of every network in the thesis: it keeps for each parameter a running average
of the gradient, which smooths the noise of the mini-batches, and a running average of its
square, by whose root the step is divided, so that every parameter moves by about $\eta$
whatever the scale of its gradient; \emph{weight decay} shrinks every parameter towards
zero by a small fraction at every step and is the simplest guard against overfitting
(\cref{sec:bg:learning:generalisation}).  The learning rate is the one quantity set by
hand: the stateful model uses $\eta = 10^{-4}$ throughout, the tabular networks a
\emph{one-cycle} schedule that raises $\eta$ from a small value to $10^{-3}$ and lowers
it again over the epochs, so that the early steps are large and the last ones fine.

\subsection{Backpropagation}
\label{sec:bg:learning:backprop}

The gradient $\nabla_{\boldsymbol{\theta}} L$ has one component per parameter, millions of
them, and \emph{backpropagation} \cite{rumelhart1986backprop} computes all of them at the
cost of about two forward passes.  A network is a chain of simple operations, each with a
known derivative.  The forward pass computes every intermediate quantity from the input to
the loss and keeps it; the backward pass applies the chain rule in reverse order, starting
from the derivative of the loss with respect to the logits, \cref{eq:xent-grad}, and
multiplying at each step back by the local derivative of the operation.  Each step reuses
the result of the step before, and the derivative with respect to a weight matrix is the
outer product of the gradient arriving at its output and the input it multiplied, which is
why the forward quantities are kept.  A network with branches, such as the stateful model
of \cref{ch:models}, in which the same state feeds several modules, sums the gradients
that arrive along the different branches; a recurrent network applies the same rule
through time (\cref{sec:bg:tbptt}).  In PyTorch \cite{paszke2019pytorch} every operation
records its inputs as it runs and the backward pass is generated automatically for any
composition of operations (\emph{automatic differentiation}), so no derivative is written
by hand.  The computation of each model is given with the model, in
\cref{sec:bg:sequences,sec:bg:graphs,sec:bg:tabular} and in \cref{ch:models}.

\paragraph{Activation functions and normalisation.}  The activation $\sigma$ of a
layer (\cref{sec:bg:mlp}) is, in the networks of the thesis, the SiLU, $\operatorname{SiLU}(a) = a /
(1 + e^{-a})$ \cite{elfwing2018silu,hendrycks2016gelu}, a smooth version of the rectifier
$\max(0, a)$ that passes positive inputs almost unchanged and suppresses negative ones;
the logistic sigmoid $\sigma(a) = 1 / (1 + e^{-a})$, which maps any number into $(0, 1)$,
and the hyperbolic tangent, which maps it into $(-1, 1)$, are the activations of the gates
of recurrent units (\cref{sec:bg:rnn}).  \emph{Layer normalisation} \cite{ba2016layernorm}
rescales a vector to zero mean and unit variance across its components and then applies a
learned scale and shift, $\operatorname{LN}(\mathbf{h}) = \boldsymbol{\gamma} \odot
(\mathbf{h} - \mu) / \sqrt{s^2 + \epsilon} + \boldsymbol{\beta}$ with $\mu$ and $s^2$ the
mean and variance of the components of $\mathbf{h}$; it keeps the quantities that flow
through a deep or recurrent network on a fixed scale, which keeps the gradients
from vanishing or exploding as they are multiplied back through many layers.

\subsection{Class imbalance: weights, the focal loss and the decision rule}
\label{sec:bg:learning:imbalance}

The cross-entropy treats every sample alike, so on data in which one class is far more
frequent than the others (89\% \textsc{const} at one minute) the model learns first of all
to predict that class, because that is where most of the loss is.  Three corrections are
used in the thesis.  \emph{Class weights} multiply the loss of each sample by a weight
$\alpha_{y}$ that depends on its class, $\ell = -\alpha_{y} \log \hat{p}_{y}$, larger for
the rare classes (2 for \textsc{down} and \textsc{up}, 1 for \textsc{const} in the stateful
model; the inverse of the class frequency in the trees and the tabular networks), so that
the rare classes count as much as the frequent one in the gradient.  The \emph{focal loss}
\cite{lin2017focal},
\begin{equation}
  \ell_{\mathrm{focal}}(\hat{\mathbf{p}}, y) = -\alpha_{y} (1 - \hat{p}_{y})^{\gamma} \log \hat{p}_{y},
  \qquad \gamma = 2,
  \label{eq:focal}
\end{equation}
multiplies the cross-entropy by $(1 - \hat{p}_y)^{\gamma}$, a factor that is close to one
when the model is wrong and close to zero when it is already right: a sample that the model
classifies with probability $0.9$ contributes one hundredth of what it would under the
cross-entropy, so training concentrates on the samples it gets wrong, which under
imbalance are mostly those of the rare classes.  Its gradient with respect to the logits is
the cross-entropy gradient \cref{eq:xent-grad} scaled by the same factor and corrected by
a term that vanishes as $\hat{p}_y \to 1$, so backpropagation proceeds unchanged.

The third correction acts on the \emph{decision}, not on the training.  A classifier is
scored by which class it predicts, $\hat{y} = \arg\max_c \hat{p}_c$, and the class that is
most probable is not the one that maximises balanced accuracy.  Balanced accuracy weighs
every class equally, which is the same as scoring the classifier on a population in which
the three classes are equally frequent; on such a population the posterior probability of
class $c$ is, by Bayes' rule, proportional to $p(y = c \mid \mathbf{x}) / \pi_c$, where
$\pi_c$ is the frequency of the class in the data the model learned from.  The decision
that maximises balanced accuracy is therefore
\begin{equation}
  \hat{y} = \arg\max_c \frac{\hat{p}_c}{\pi_c} = \arg\max_c \bigl( z_c - \log \pi_c \bigr),
  \label{eq:logit-adjustment}
\end{equation}
the most probable class \emph{after} a constant has been subtracted from the logit of each
class, larger for the frequent classes \cite{menon2021logitadjustment}.  The constants are
exact when the probabilities $\hat{\mathbf{p}}$ are correct; when they are not, and the
focal loss and the class weights already distort them, the constants are better fitted than
computed.  The calibration rule of \cref{sec:models:protocol} is exactly this: two
constants added to the \textsc{down} and \textsc{up} logits of every model,
$\boldsymbol{\delta}$ in \cref{sec:bg:notation}, chosen on the validation day to maximise
balanced accuracy and applied unchanged afterwards, so that every model is compared with
its decision rule placed the same way.

\subsection{Generalisation, overfitting and early stopping}
\label{sec:bg:learning:generalisation}

A model with millions of parameters can reach a very small loss on the training samples
by memorising them, and such a model is useless on new data; what matters is the loss, or
the balanced accuracy, on samples the model has not seen, its \emph{generalisation}.  The
thesis guards against \emph{overfitting} in the standard ways \cite{goodfellow2016deep}.
The data are split by date into training, validation and test periods
(\cref{sec:data:splits}): the parameters are fitted on the training days, every choice
about the model is made on the validation day, and the test days are read once at the end.
\emph{Early stopping} scores the model on a held-out day after every epoch (12 March for
the trees and the tabular networks) and keeps the parameters of the best epoch, stopping
when the score has not improved for a fixed number of epochs; it is the main regulariser
of the trees and of the tabular networks.  \emph{Weight decay}, the shrinkage
term of AdamW (\cref{sec:bg:learning:sgd}), keeps the weights small.  \emph{Dropout} \cite{srivastava2014dropout}
sets a random fraction of the hidden units to zero at every training step, so that no unit
can rely on a specific other unit being present, and uses all of them, scaled, at
prediction time; the tabular networks use a rate of $0.15$, the predictor of the stateful
model $0.1$.  And the \emph{deduplicated} evaluation of \cref{sec:data:labels}, which
keeps one sample per market per fifteen minutes, guards against a subtler form of the
same problem: ten million samples that are near-copies of one another make a model look
far better on the busy markets than it is on independent ones.
